\documentclass[11pt]{article}
\usepackage[utf8]{inputenc}
\usepackage[T1]{fontenc}
\usepackage[a4paper,margin=1in]{geometry}
\usepackage{lmodern}
\usepackage{microtype}
\usepackage{amsmath,amssymb,mathtools,amsthm}
\usepackage{mathrsfs}
\usepackage{enumitem}
\usepackage{tikz-cd}
\usepackage[hidelinks,hypertexnames=false]{hyperref}
\setlength{\parindent}{0pt}
\setlength{\parskip}{0.55em}
\emergencystretch=3em
\hbadness=10000
\providecommand{\Spec}{\operatorname{Spec}}
\providecommand{\Pic}{\operatorname{Pic}}
\providecommand{\Lie}{\operatorname{Lie}}
\providecommand{\Hom}{\operatorname{Hom}}
\providecommand{\Ker}{\operatorname{Ker}}
\providecommand{\Coker}{\operatorname{Coker}}
\providecommand{\rank}{\operatorname{rank}}
\providecommand{\Div}{\operatorname{Div}}
\providecommand{\Ext}{\operatorname{Ext}}
\providecommand{\End}{\operatorname{End}}
\providecommand{\degop}{\operatorname{deg}}
\providecommand{\im}{\operatorname{im}}
\providecommand{\cO}{\mathcal O}
\providecommand{\Z}{\mathbb Z}
\providecommand{\Q}{\mathbb Q}
\providecommand{\R}{\mathbb R}
\providecommand{\C}{\mathbb C}
\providecommand{\Gm}{\mathbb G_m}
\providecommand{\Ga}{\mathbb G_a}
\providecommand{\Ql}{\mathbb Q_\ell}
\providecommand{\Zl}{\mathbb Z_\ell}
\providecommand{\Fp}{\mathbb F_p}
\providecommand{\et}{\mathrm{et}}
\providecommand{\etale}{\mathrm{\acute et}}
\providecommand{\op}{\mathrm{op}}
\providecommand{\id}{\mathrm{id}}
\providecommand{\Fil}{\operatorname{Fil}}
\providecommand{\gr}{\operatorname{gr}}
\providecommand{\pr}{\operatorname{pr}}
\providecommand{\Biext}{\operatorname{Biext}}
\providecommand{\BT}{\operatorname{BT}}
\providecommand{\cE}{\mathcal E}
\providecommand{\cF}{\mathcal F}
\providecommand{\cG}{\mathcal G}
\providecommand{\cH}{\mathcal H}
\providecommand{\cM}{\mathcal M}
\providecommand{\cR}{\mathcal R}
\providecommand{\cT}{\mathcal T}
\providecommand{\cL}{\mathcal L}
\providecommand{\cD}{\mathcal D}
\providecommand{\Ab}{\operatorname{Ab}}
\providecommand{\Sch}{\operatorname{Sch}}
\providecommand{\an}{\mathrm{an}}
\providecommand{\red}{\mathrm{red}}
\providecommand{\nil}{\mathrm{nil}}
\providecommand{\can}{\mathrm{can}}
\providecommand{\ord}{\mathrm{ord}}
\DeclareMathOperator{\Tor}{Tor}
\newenvironment{statement}[1]{\par\medskip\noindent\textbf{#1.}\ }{\par\medskip}
\newenvironment{prooftext}{\par\noindent\emph{Proof.}\ }{\par\medskip}
\newenvironment{remarktext}[1]{\par\medskip\noindent\textbf{#1.}\ }{\par\medskip}
\begin{document}
\begin{center}
{\Large SGA 7-I: Expos\'e IX, continuation}\\[0.4em]
{\large N\'eron models, the Picard--Lefschetz formula, and analytic comparison}
\end{center}

\section*{12. N\'eron model of a Jacobian and the Picard--Lefschetz formula}

We first recall the following general results of M. Raynaud.  They are used in the form needed below, where the base is a strictly local trait.

\begin{statement}{Theorem 12.1 (M. Raynaud)}
Let \(S\) be a strictly local trait and let \(f:X\to S\) be proper and flat, with fibres of dimension one, such that
\begin{equation}\tag{12.1.1}
        H^1(X,\cO_X)=0
\end{equation}
and such that the local rings of \(X\) are factorial.  For each maximal point \(x_i\) of the special fibre, put
\begin{equation}\tag{12.1.2}
        a_i=\mu_i\,[k(x_i):k]_{\rm rad},
\end{equation}
where \(\mu_i\) is the multiplicity of the component through \(x_i\), and where the radical degree is a power of the characteristic exponent of the residue field.  Let \(d\), respectively \(d^0\), be the greatest common divisor of the integers \(a_i\), respectively of their separable parts.  Thus \(d=d^0\) when the residue field is perfect.  Put
\begin{equation}\tag{12.1.3}
        P=\Pic_{X/S}\colon (\Sch)/S^{\op}\longrightarrow (\Ab).
\end{equation}
Consider the following conditions:
\begin{enumerate}[label=\arabic*)]
\item \(d=1\).
\item \(X\) has a zero-cycle of degree \(1\).
\item There exists an invertible module \(L\) on \(X\times P^0\) whose corresponding homomorphism \(P^0\to \Pic^0_{X/S}=P^0\) is the canonical inclusion.
\item \(X\) is cohomologically flat over \(S\); equivalently, it is cohomologically flat in dimension zero, and the canonical map
\[
        k\longrightarrow H^0(X_s,\cO_{X_s})
\]
is an isomorphism, with \(d=1\).
\item \(P\) is representable, equivalently by a smooth group scheme over \(S\).
\item \(P^0\) is representable, equivalently by a smooth group scheme over \(S\).
\item \(P^0\) is separated over \(S\), i.e. satisfies the valuative criterion of separatedness; equivalently, every section of \(P^0\) over \(S\) which is zero on the generic point is zero.
\item \(d^0=1\).
\end{enumerate}
Then one has the chain of implications displayed in Raynaud's theorem; in particular, when \(d=d^0\), for example when \(k\) is perfect, the eight conditions are equivalent.

Assume now that \(d^0=1\), so that \(P\) is represented by a smooth group scheme over \(S\).  Then the schematic closure \(E\) of the unit section of \(P\) is an \(S\)-\'etale subgroup scheme.  If \(I\) denotes the set of reduced irreducible components \(X_i\) of the special fibre, one obtains an isomorphism
\begin{equation}\tag{12.1.5}
        \Z^I/\Z\xrightarrow{\sim}E
\end{equation}
by associating to \(i\in I\) the section of \(P\) defined by the invertible sheaf \(\cO_X(X_i)\).  Moreover
\begin{equation}\tag{12.1.6}
        P^0\cap E=\text{the zero section of }P.
\end{equation}
The quotient
\begin{equation}\tag{12.1.7}
        Q=P/E
\end{equation}
is represented by a smooth separated group scheme over \(S\), and is N\'eronian in the sense that it satisfies the N\'eron mapping property on its special fibre.  If, in addition, the geometric generic fibre of \(X\) is integral, there is a canonical exact sequence
\begin{equation}\tag{12.1.8}
        0\longrightarrow A\longrightarrow Q\xrightarrow{\degop}\Z_S\longrightarrow 0,
\end{equation}
where \(A\) is smooth, separated, and of finite type over \(S\), and is the N\'eron model of
\begin{equation}\tag{12.1.9}
        A_\eta=\Pic^0_{X_\eta/\eta}.
\end{equation}
The composite \(P^0\to P\to Q\) induces an isomorphism
\begin{equation}\tag{12.1.10}
        P^0\xrightarrow{\sim}A^0.
\end{equation}
\end{statement}

The case of interest for us is the following.  Suppose \(X_\eta\) is smooth, so that, in view of (12.1.1), the geometric generic fibre is smooth and connected, and suppose that the special fibre has a zero-cycle of degree \(1\).  Then all the preceding conditions hold.  The group \(A_\eta=\Pic^0_{X_\eta/\eta}\) is an abelian scheme on the generic point, and (12.1.8) gives a very explicit construction of its N\'eron model.  In particular (12.1.10) gives
\begin{equation}\tag{12.1.12}
        A^0_s\simeq \Pic^0_{X_s/s}.
\end{equation}
It follows that \(A_\eta\) has semi-stable reduction exactly when the maximal torus in \(\Pic^0_{X_s/s}\) has the expected semi-stable form, equivalently in the regular case when the only singularities of the geometric special fibre are ordinary quadratic points.  This last verification follows from the usual description of the Picard scheme of a proper separable curve over a field and from EGA IV, 21.8.5; it is also an immediate consequence of the geometric monodromy theorem together with the Galois criterion for semi-stable reduction.

Raynaud's proof is cited from the papers used above.  We shall need the theorem only in the special situation just described, and from now on we restrict to the case in which the singular points of the special fibre are ordinary quadratic points.

\subsection*{12.3. The toric part of the Picard scheme of a singular curve}

Let \(C\) be a proper separable curve over a separably closed field \(k\), and let \(\widetilde C\) be its normalization.  The Picard scheme \(\Pic^0_{\widetilde C/k}\) has reductive rank zero, hence the maximal torus of \(\Pic^0_{C/k}\) is contained in
\[
        N=\Ker\bigl(\Pic^0_{C/k}\longrightarrow \Pic^0_{\widetilde C/k}\bigr).
\]
By EGA IV, 21.8.5, the group \(N\) is an extension of a torus \(T\), described below, by a connected unipotent group.  Thus \(T\) is canonically the maximal torus of \(\Pic^0_{C/k}\).

Let \(J\) be the set of irreducible components of \(\widetilde C\).  For \(x\in C\), let \(J(x)\) be the set of branches of \(C\) through \(x\), i.e. the set of points of \(\widetilde C\) above \(x\).  Put
\begin{equation}\tag{12.3.1}
        R(x)=\Z^{J(x)}/\Z,
\end{equation}
where \(\Z\to\Z^{J(x)}\) is the diagonal map, and
\begin{equation}\tag{12.3.2}
        R=\bigoplus_x R(x),
\end{equation}
the sum being taken only over points with at least two branches.  Then
\begin{equation}\tag{12.3.3}
        T=\Hom(R,\Gm).
\end{equation}
Equivalently, the character group \(M\) of \(T\) is
\begin{equation}\tag{12.3.4}
        M=\Hom(T,\Gm)=\Ker\left(\bigoplus_x\Z^{J(x)}\longrightarrow \Z^{J}\right),
\end{equation}
where the map is induced by the natural map sending a branch to the component on which it lies.

We shall use only the case in which at most two branches pass through a point.  If \(I\) is the set of points with two branches, then, after choosing an order of the two branches at each such point,
\begin{equation}\tag{12.3.6}
        R\simeq \Z^I,
        \qquad
        M=\Ker\bigl(\Z^I\longrightarrow \Z^J\bigr).
\end{equation}
Equivalently one can describe \(M\) by the graph \(\Gamma(C)\) whose vertices are the irreducible components and whose edges are the unordered pairs of branches meeting at a singular point.  Then
\begin{equation}\tag{12.3.7}
        M\simeq H_1(\Gamma(C),\Z),
        \qquad
        R\simeq H^1(\Gamma(C),\Z).
\end{equation}
The canonical homomorphism from the branch description to the graph description is induced by the map from the set of branches over the double points to the set of components.

\subsection*{12.4. The pairing attached to ordinary double points}

We now place ourselves in the situation fixed at the end of 12.2.  By (12.1.12) and by the preceding description applied to \(C=X_s\), the character group \(M\) of the maximal torus \(T_s\) of the special fibre \(A^0_s\) of the N\'eron model is
\begin{equation}\tag{12.4.1}
        M=\Ker\bigl(\Z^I\longrightarrow \Z^J\bigr),
\end{equation}
where \(I\) is the set of singular points of \(X_s\), and \(J\) the set of irreducible components.  The map \(\Z^I\to\Z^J\) is obtained by assigning, to a double point, the difference of the two components which contain the two branches through that point.  Thus, after choosing an order \((C'_i,C''_i)\) of the two branches through \(x_i\), one has
\begin{equation}\tag{12.4.2}
        d(x_i)=C'_i-C''_i.
\end{equation}
The operator \(d\), and the isomorphism (12.4.1), depend on these choices.  However, if
\begin{equation}\tag{12.4.3}
        \varphi_{x_i}:M\longrightarrow \Z
\end{equation}
is the composite of \(M\to\Z^I\) with the projection onto the factor corresponding to \(x_i\), then \(\varphi_{x_i}\) is determined up to sign, and the sign changes exactly when the order of the two branches is reversed.  Consequently the bilinear form
\begin{equation}\tag{12.4.4}
        \varphi_{x_i}\otimes\varphi_{x_i}\in
        M^\vee\otimes_\Z M^\vee\simeq \Hom(M\otimes M,\Z)
\end{equation}
is canonically determined by the singular point \(x_i\).  We may therefore put
\begin{equation}\tag{12.4.5}
        u=\sum_i \varphi_{x_i}\otimes\varphi_{x_i}:
        M\otimes_\Z M\longrightarrow \Z.
\end{equation}
This pairing is well-defined, independent of all choices of signs.  It is positive definite, being the restriction to \(M\subset\Z^I\) of the standard quadratic form \(\sum_i X_i^2\) on \(\Z^I\).

The notation \(M\) for the character group of \(T_s\) differs from the notation used earlier, where \(M'\) denoted the corresponding group for the dual abelian variety.  Since here \(A_\eta\) is a Jacobian, it has its canonical principal polarization, giving a canonical isomorphism
\begin{equation}\tag{12.4.6}
        A_\eta\xrightarrow{\sim}A'_\eta.
\end{equation}
We shall use this to identify \(A_\eta\) and \(A'_\eta\), and hence \(M\) and \(M'\).  For every prime \(\ell\), the monodromy pairing may therefore be viewed as a pairing
\begin{equation}\tag{12.4.7}
        u_\ell:M\otimes\Z_\ell\times M\otimes\Z_\ell\longrightarrow \Z_\ell.
\end{equation}

\begin{statement}{Theorem 12.5 (Picard--Lefschetz formula)}
Let \(S\) be a strictly local trait, and let \(X\) be regular, projective and flat over \(S\), with fibres of dimension one.  Assume that the geometric generic fibre is smooth and connected, and that the geometric special fibre has no singularities other than ordinary quadratic points.  With the notation introduced in 12.4, for every prime \(\ell\) the monodromy pairing (12.4.7) is the scalar extension to \(\Z_\ell\) of the integral pairing (12.4.5).
\end{statement}

This proves the monodromy formula used earlier in the case of a Jacobian equipped with its canonical polarization.  By the reduction argument already given in the preceding section, this is enough to imply the general statement for abelian varieties.

For \(\ell\ne p\), the proof is essentially the cohomological Picard--Lefschetz formula for
\begin{equation}\tag{12.7.1}
        H^1(X_{\bar\eta},\Z_\ell)
        \simeq \Hom(T_\ell(A_\eta),\Z_\ell).
\end{equation}
This formula is proved later in the seminar in a more general situation; in the present case it is exactly the assertion above, once the definition of the monodromy pairing is unfolded.  It remains to treat the case \(\ell=p\).  We reduce it to characteristic zero.

The reduction uses the same principle as the proof of the essential tameness theorem for monodromy on \(T_\ell\).  We may suppose \(S\) complete.  Let \(W\) be a Cohen \(p\)-ring with residue field \(k\), and choose a homomorphism \(W\to V\) inducing the identity on residue fields.  Let \(L\) be the set of singular points of \(X_s\).  Consider the formal moduli space \(\mathfrak S\) of \(X_s\) over \(W\), with universal curve
\[
        \mathfrak X\longrightarrow\mathfrak S.
\]
As in the local deformation theory recalled earlier, \(\mathfrak S\) is isomorphic to the spectrum of a formal power series ring \(W[[T_\lambda]]_{\lambda\in L}\), and in \(\mathfrak S\) there is a normal-crossings divisor
\begin{equation}\tag{12.8.1}
        D=\bigcup_{\lambda\in L}D_\lambda,
\end{equation}
where the component \(D_\lambda\) corresponds to smoothing all singularities except the point indexed by \(\lambda\).  The given family over \(S\) is obtained by pullback from the universal family:
\begin{equation}\tag{12.8.2}
\begin{tikzcd}
X\arrow[r]\arrow[d]&\mathfrak X\arrow[d]\\
S\arrow[r,"g"]&\mathfrak S.
\end{tikzcd}
\end{equation}
Let
\begin{equation}\tag{12.8.3}
        P=\Pic^0_{\mathfrak X/\mathfrak S}.
\end{equation}
Since the special fibre is geometrically reduced of dimension one, \(\mathfrak X\) is cohomologically flat in dimension zero, and \(P\) is formally smooth over \(\mathfrak S\).  We use the representability of this Picard functor by a smooth group object over \(\mathfrak S\).  Alternatively, Artin's representability theorem supplies an algebraic space, which suffices for the argument.

Let
\begin{equation}\tag{12.8.4}
        U=\mathfrak S-\operatorname{Supp}D.
\end{equation}
Over \(U\), the universal curve is smooth and \(P_U=\Pic^0_{\mathfrak X_U/U}\) is an abelian scheme.  Pulling back along \(g:S\to U\) gives the abelian variety \(A_\eta\) we are studying.  The isomorphism (12.1.10) gives, on \(\mathfrak S\),
\begin{equation}\tag{12.8.5}
        P^0\times_{\mathfrak S} S\simeq A^0.
\end{equation}
Since the special fibre of \(P\) is \(\ell\)-divisible for every \(\ell\), so is \(P\) itself.  We may therefore consider \(T_\ell(P)\) over \(\mathfrak S\); its pullback to \(S\) is \(T_\ell(A_\eta)\).  Applying the constructions of the previous sections gives a filtration by the fixed part and the toric part
\begin{equation}\tag{12.8.6}
        T_\ell(P)^f\subset T_\ell(P)^t\subset T_\ell(P).
\end{equation}
Over \(U\) the canonical autoduality of the Jacobian gives a pairing
\begin{equation}\tag{12.8.7}
        T_\ell(P_U)\times T_\ell(P_U)\to\Z_\ell(1).
\end{equation}
For this pairing, the fixed part and the toric part are mutual orthogonals.  Thus the quotient \(E=T_\ell(P_U)^t/T_\ell(P_U)^f\) is canonically dual to \(F=T_\ell(P_U)/T_\ell(P_U)^t\), and extends canonically over \(\mathfrak S\) as
\begin{equation}\tag{12.8.8}
        E=D(F),\qquad F=T_\ell(P)^t/T_\ell(P)^f.
\end{equation}
We are therefore in the typical situation of a mixed extension.  The group \(T_\ell(P_U)\) is a mixed extension of \(E\) by \(F\), where \(E\) and \(F\) arise from two dual extensions.  Schematically the situation is
\begin{equation}\tag{12.8.9}
\begin{tikzcd}[column sep=large]
0\arrow[r]&F\arrow[r]&R\arrow[r]&M\arrow[r]&0\\
&&T_\ell(P_U)\arrow[u]\arrow[d]&&\\
0\arrow[r]&D(M)\arrow[r]&R^\vee\arrow[r]&D(F)\arrow[r]&0.
\end{tikzcd}
\end{equation}
The only difference from the earlier discussion is that the resulting pairing takes values not in a canonically fixed copy of \(\Z_\ell\), but in
\[
        \Gamma(U,\cO_U^*)/\Gamma(\mathfrak S,\cO_{\mathfrak S}^*)\otimes\Z_\ell,
\]
which is identified with \(\Z^L\otimes\Z_\ell\) by the functions defining the components \(D_\lambda\).  This gives a canonical pairing
\begin{equation}\tag{12.8.11}
        u: M\otimes M\longrightarrow\Z^L.
\end{equation}
The monodromy pairing for the given curve over \(S\) is the composite of this pairing with the homomorphism \(\Z^L\to\Z\) induced by \(g:S\to\mathfrak S\).  The latter homomorphism is given by an element
\begin{equation}\tag{12.8.12}
        (d_\lambda)_{\lambda\in L}\in\Z^L,
        \qquad d_\lambda\geq 1,
\end{equation}
where \(d_\lambda\) is the multiplicity with which the pullback of \(D_\lambda\) appears on \(S\).  For fixed \(\lambda\), one has \(d_\lambda=1\) precisely when \(X\) is regular at the corresponding double point.  Since regularity is assumed in Theorem 12.5, all \(d_\lambda\)'s are equal to \(1\).  It remains to prove the following local assertion.

\begin{statement}{Lemma 12.8.13}
For \(\lambda\in L\), the composite of (12.8.11) with the projection \(\Z^L\to\Z\) onto the factor \(\lambda\) is the form \(\varphi_{x_\lambda}\otimes\varphi_{x_\lambda}\) of (12.4.4).
\end{statement}

Complete the strictly local ring at the generic point of \(D_\lambda\) and pull the family back to this trait.  The new special fibre has exactly one ordinary quadratic singularity.  The preceding constructions commute with this localization: the torus, its character group, the monodromy form, and the specialization morphism all localize compatibly.  Thus the assertion reduces to Theorem 12.5 in the case where the base is a trait of characteristic zero and the special fibre has a single ordinary double point.  This is the classical Picard--Lefschetz case.  The reduction proves the theorem.

\begin{remarktext}{Remarks 12.9--12.10}
The reduction just made also reduces to the classical situation in which the special fibre has one double point.  In that case there are one or two irreducible components; these correspond respectively to \(\rank M=1\) and \(\rank M=0\), the latter being the case of good reduction for the Jacobian.  Without assuming regularity of \(X\), the same argument gives a more general formula in which the contribution of each double point is multiplied by the integer \(d_\lambda\) of (12.8.12): the monodromy pairing is obtained from
\begin{equation}\tag{12.10.1}
        \sum_{\lambda\in L}d_\lambda\,
        \varphi_{x_\lambda}\otimes\varphi_{x_\lambda}.
\end{equation}
\end{remarktext}

\section*{13. Links with transcendental theory: the complex analytic case}

For general facts on analytic spaces we refer to the standard sources.  We review only the analytic dictionary needed to compare the preceding algebraic constructions with Hodge-theoretic ones.

\subsection*{13.1. Analytic groups over a base}

Let \(S\) be a complex analytic space, and let \(A\) be a commutative analytic group over \(S\).  Assume \(A\) is smooth over \(S\), and put
\begin{equation}\tag{13.1.1}
        E=\Lie(A/S),
\end{equation}
the locally free \(\cO_S\)-module pulled back from the relative tangent sheaf along the zero section.  We shall use the exponential homomorphism
\begin{equation}\tag{13.1.2}
        \exp:E\longrightarrow A,
\end{equation}
where \(E\) denotes the analytic vector bundle with sheaf of holomorphic sections \(E\).  It is defined locally near the zero section by the standard theory of formal groups in characteristic zero and then extended by multiplicativity.  Its image is the open analytic subgroup \(A^0\), the union of the neutral connected components of the fibres of \(A\).  If \(R=\Ker(\exp)\), then one has an exact sequence of relative analytic groups
\begin{equation}\tag{13.1.3}
        0\longrightarrow R\longrightarrow E\longrightarrow A^0\longrightarrow 0.
\end{equation}
The group \(R\) is \(S\)-\'etale and is a closed analytic subspace of \(E\) precisely when \(A^0\) is separated over \(S\).

The exact sequence (13.1.3) is functorial in \(A\) and commutes with arbitrary base change.  In this way one obtains a base-change-compatible equivalence between smooth relative analytic groups with connected fibres and pairs \((E,R)\), where \(E\) is a locally free analytic vector bundle on \(S\) and \(R\subset E\) is an \(S\)-\'etale analytic subgroup.

When \(A\) is abeloid over \(S\), meaning smooth, proper and separated over \(S\) with connected fibres, then \(A=A^0\), and \(R\) is locally constant.  Its rank in a fibre is twice the complex rank of \(E\).  Conversely, an abeloid analytic group is described by such a local system \(R\) together with a complex analytic vector-bundle structure on the real analytic bundle \(R\otimes_\Z\R\).

\subsection*{13.2. Formal completion along a fibre}

Fix a point \(s\in S\).  For \(n\geq 0\), let \(S_n\) be the \(n\)-th infinitesimal neighbourhood of \(s\) in \(S\), and for an analytic space \(X\) over \(S\), put \(X_n=X\times_S S_n\).  The system \(\widehat X=(X_n)_n\) is the formal completion of \(X\) along the fibre \(X_s\).  It is functorial in \(X\), and formation of \(\widehat X\) commutes with finite projective limits.

\begin{statement}{Lemma 13.2.1}
Let \(A\) be as in 13.1, and suppose that:
\begin{enumerate}[label=\alph*)]
\item \(A=A^0\);
\item \(R\) is constant in a neighbourhood of \(s\);
\item \(R_s\) generates the complex vector space \(E_s\).
\end{enumerate}
If \(A'\) is another smooth analytic group over \(S\), then the natural map from germs of homomorphisms \(A\to A'\) near \(s\) to homomorphisms of formal completions
\begin{equation}\tag{13.2.1.1}
        \varinjlim_{U\ni s}\Hom_U(A_U,A'_U)
        \longrightarrow
        \Hom_{\widehat S}(\widehat A,\widehat A')
\end{equation}
is bijective.
\end{statement}

\begin{prooftext}
Using (13.1.3), we may replace a homomorphism \(A\to A'\) by a homomorphism of the corresponding vector bundles carrying \(R\) into \(R'\).  Passing to formal completions amounts to passing from an \(\cO_s\)-linear homomorphism to its completion.  The injectivity is clear.  For surjectivity, Nakayama's lemma and the hypothesis that \(R_s\) generates \(E_s\) show that any compatible formal homomorphism of vector bundles sends the local system \(R\) into \(R'\) in a neighbourhood of \(s\).  Thus it comes from a genuine germ of homomorphisms.
\end{prooftext}

Consequently, under the hypotheses of the lemma, the analytic group \(A\) is determined, up to unique isomorphism in a neighbourhood of \(s\), by its formal completion \(\widehat A\).

\subsection*{13.3. The analytic Raynaud extension}

Let \(A\) be as in 13.1 with \(A=A^0\).  Write \(R_s\) for the lattice in the fibre.  After shrinking \(S\) around \(s\), the group of germs of sections of \(R\) gives an injective homomorphism from the constant group \(R_s\) to \(R\).  Its image is the largest constant subgroup of \(R\); call it the fixed part:
\begin{equation}\tag{13.3.2}
        R^f\subset R,
        \qquad R^f_s\xrightarrow{\sim}R_s^f.
\end{equation}
The quotient
\begin{equation}\tag{13.3.3}
        A^\natural=E/R^f
\end{equation}
fits into an exact sequence
\begin{equation}\tag{13.3.4}
        0\longrightarrow R/R^f\longrightarrow A^\natural\longrightarrow A\longrightarrow 0.
\end{equation}
The morphism \(A^\natural\to A\) is \(\etale\), hence induces an isomorphism
\begin{equation}\tag{13.3.5}
        \Lie(A^\natural/S)\xrightarrow{\sim}\Lie(A/S)=E.
\end{equation}
For every \(n\geq 0\), it also induces an isomorphism on the \(n\)-th infinitesimal neighbourhood of the fibre, and therefore
\begin{equation}\tag{13.3.6}
        \widehat{A^\natural}\xrightarrow{\sim}\widehat A.
\end{equation}
Thus \(A^\natural\) is the analytic counterpart of the group considered algebraically in the Raynaud construction.

Suppose now that a complex torus
\begin{equation}\tag{13.3.7}
        T_s\subset A_s
\end{equation}
is given.  Let \(T\) be the constant torus over \(S\) with fibre \(T_s\).  The induced homomorphism of formal completions comes, by Lemma 13.2.1, from a unique germ of analytic homomorphisms
\begin{equation}\tag{13.3.9}
        T\longrightarrow A^\natural.
\end{equation}
After shrinking \(S\), this homomorphism is defined on all of \(S\).  Its differential identifies \(\Lie(T/S)\) with a locally direct-factor submodule \(E^t\subset E\):
\begin{equation}\tag{13.3.10}
        \Lie(T/S)\simeq E^t\subset E.
\end{equation}
Let
\begin{equation}\tag{13.3.11}
        R^t\subset R_s,
        \qquad R^t\subset R,
\end{equation}
be the sublattice defined by the torus \(T_s\) and the corresponding constant subgroup.  The immersion condition implies, after shrinking \(S\), that \(T\to A^\natural\) is a closed immersion.  Put
\begin{equation}\tag{13.3.12}
        B=A^\natural/T.
\end{equation}
Then \(A^\natural\) is an extension
\begin{equation}\tag{13.3.13}
        0\longrightarrow T\longrightarrow A^\natural\longrightarrow B\longrightarrow 0,
\end{equation}
where \(B\) is smooth and separated over \(S\), with connected fibres, and
\begin{equation}\tag{13.3.14}
        \Lie(T)=E^t.
\end{equation}
The exponential sequence for \(B\) is
\begin{equation}\tag{13.3.15}
        0\longrightarrow R/R^t\longrightarrow E/E^t\longrightarrow B\longrightarrow 0.
\end{equation}
Since
\begin{equation}\tag{13.3.16}
        R^t/R^f=(R/R^f)^t
\end{equation}
is constant, \(B\) is abeloid near \(s\) precisely when its fibre \(B_s=A_s/T_s\) is abeloid.  Therefore every expression of \(A_s\) as an extension of an abeloid analytic group by a complex torus gives, near \(s\), a unique extension structure of \(A^\natural\) by a torus whose fibre is the given torus.

\subsection*{13.4. Mixed Hodge structures}

We reformulate the dictionary (13.1.4).  Let \(\cR\) be the sheaf of sections of \(R\), and put
\begin{equation}\tag{13.4.1}
        \cE(R)=\cR\otimes_\Z\cO_S.
\end{equation}
When \(R\) is locally constant this is a locally free \(\cO_S\)-module.  The inclusion \(R\subset E\) gives a canonical homomorphism
\begin{equation}\tag{13.4.2}
        \cE(R)\longrightarrow E.
\end{equation}
Let \(F\) be its kernel, so that
\begin{equation}\tag{13.4.3}
        0\longrightarrow F\longrightarrow \cE(R)\longrightarrow E
\end{equation}
is exact on the left.  In the most useful case the last map is surjective; then
\begin{equation}\tag{13.4.4}
        0\longrightarrow F\longrightarrow \cE(R)\longrightarrow E\longrightarrow0
\end{equation}
is exact.  The pair \((E,R)\) is then equivalent to the pair \((R,F)\), where \(R\) is locally constant with finitely generated free \(\Z\)-modules as fibres and \(F\subset \cE(R)\) is a locally direct-factor submodule satisfying
\begin{equation}\tag{13.4.5}
        F\cap\overline F=0
\end{equation}
in the usual fibrewise sense.  We regard \(\cE(R)\) as endowed with the Hodge filtration
\begin{equation}\tag{13.4.6}
        \Fil^0\cE(R)=\cE(R),\qquad
        \Fil^1\cE(R)=F,\qquad
        \Fil^2\cE(R)=0.
\end{equation}
The exact sequence (13.4.3) is functorial in \(A\) and commutes with base change.

Returning to 13.3, set
\begin{equation}\tag{13.4.7}
        \cE^f=R^f\otimes\cO_S,
        \qquad
        \cE^t=R^t\otimes\cO_S.
\end{equation}
These form the increasing weight filtration
\begin{equation}\tag{13.4.8}
        \cE^f\subset\cE^t\subset\cE=\cE(R),
\end{equation}
where \(\cE^f\), \(\cE^t/\cE^f\), and \(\cE/\cE^t\) are to be thought of as of weights \(-2\), \(-1\), and \(0\), respectively.  The functoriality of (13.4.3), applied to \(T\to A^\natural\to A\), gives
\begin{equation}\tag{13.4.10}
        F(A^\natural)=F\cap \cE^f,
        \qquad
        F(T)=0.
\end{equation}
Consequently the composite \(\cE^t\to E\) induces an isomorphism
\begin{equation}\tag{13.4.11}
        \cE^t/\cE^f\xrightarrow{\sim}E^t=\Lie(T).
\end{equation}
Assume moreover that \(R_s\) generates \(E_s\), for instance in the favourable case where \(A_s\) is an extension of an abeloid analytic group by a torus.  Then the corresponding map for \(A^\natural\) is surjective, and one has exact sequences
\begin{equation}\tag{13.4.12}
        0\longrightarrow F^\natural\longrightarrow \cE^f\longrightarrow E\longrightarrow0,
\end{equation}
with
\begin{equation}\tag{13.4.13}
        F\cap \cE^t=0.
\end{equation}
For the quotient \(B=A^\natural/T\), one has
\begin{equation}\tag{13.4.14}
        0\longrightarrow F(B)\longrightarrow \cE(B)\longrightarrow \Lie(B)\longrightarrow0,
\end{equation}
and explicitly
\begin{equation}\tag{13.4.15}
        \cE(B)=\cE/\cE^t,
        \qquad
        \Lie(B)=E/E^t.
\end{equation}
Using (13.4.11) and (13.4.13), one obtains an isomorphism
\begin{equation}\tag{13.4.16}
        F^\natural\xrightarrow{\sim}F(B).
\end{equation}

\begin{statement}{Proposition 13.4.17}
In the favourable situation just described, the Hodge filtration and the weight filtration
\begin{equation}\tag{13.4.17.1}
        \cE\supset F\supset0,
        \qquad
        \cE\supset\cE^t\supset\cE^f\supset0
\end{equation}
on \(\cE=R\otimes\cO_S\) satisfy:
\begin{enumerate}[label=\alph*)]
\item the weight filtration comes by tensoring with \(\cO_S\) from a filtration
\[
        R\supset R^t\supset R^f\supset0
\]
by locally constant analytic groups whose fibres are finitely generated free \(\Z\)-modules;
\item one has
\[
        F\cap \cE^t=0,
        \qquad
        F+\cE^f=\cE;
\]
thus the filtration induced by the Hodge filtration on \(\cE^t\) is concentrated in degree \(-1\), while that induced on \(\cE/\cE^f\) is concentrated in degree \(0\);
\item the filtration induced on
\[
        \cE^t/\cE^f=(R^t/R^f)\otimes\cO_S
\]
by the Hodge filtration corresponds to the analytic abeloid group \(B\) of (13.3.13).
\end{enumerate}
\end{statement}

\begin{remarktext}{Remark 13.4.18}
In Deligne's terminology, after restriction to an open set on which \(R\) is locally constant, the preceding proposition says that \(R\), together with the Hodge and weight filtrations on \(R\otimes\cO_S\), is an analytic family of mixed Hodge structures.  This situation was one of the models for Deligne's conjectures on N\'eron theory in higher dimensions.  However, a feature present in the algebraic theory of Section 7 is not visible in this naive analytic description: when \(A\) is abeloid over the complement of a rare analytic subset, the local system \(R/R^t\) should be constant, and more precisely dual to the corresponding toric system for a suitable dual analytic group.  This follows algebraically from the duality theorem used earlier.  It is not clear from the elementary analytic methods used here whether the analogous assertion, or even unipotence of the local monodromy action, follows without additional hypotheses.  Under algebraicity hypotheses the next section shows that the expected behaviour does hold and relates it to the construction of Section 7.
\end{remarktext}

\end{document}
